% ============================================================
\section{Evaluation}
\label{sec:eval}

We ask whether a cheap geometry-aware placement routes as well as far more
expensive joint schedulers. This section describes the implementation and
benchmark suite, the baselines we compare against, and our results.

\subsection{Experimental Setup}

The method is implemented in an open-source C++20 compiler~\cite{ftqc} that parses
OpenQASM, generates the rectangular architecture and its magic-state placement,
runs the mapping of Section~\ref{sec:approach}, and routes the result with a
lattice-surgery router; the lattice dimensions are sized automatically from the
qubit count, the interaction degree, and the magic-state budget. The benchmark
suite spans \todo{$N$} circuits from three families---arithmetic
(\texttt{adder}, \texttt{multiplier}, \texttt{square\_root}), algorithmic
(\texttt{qft}, \texttt{qpe}, \texttt{bwt}, \texttt{hhl}, \texttt{factor247}), and
structured or state-preparation (\texttt{ghz}, \texttt{wstate}, \texttt{ising},
\texttt{qaoa})---with logical qubit counts from \todo{min} to \todo{max}.

We compare our (fine) Gaussian placement against two baselines. The first,
\emph{random}, scatters qubits over a sparse sublattice with $3\times3$ spacing,
which trivially satisfies safe passage (the cube condition of
Section~\ref{sec:background}), and routes it with the same router as Gaussian;
since each restart is so cheap, we run it for \todo{N} repetitions and
keep the best result, whereas the Gaussian placement is deterministic and runs
once. The second is the external WISQ compiler~\cite{wisq}. The primary metrics
are the routing-step count $|\Rout|$ (Definition~\ref{def:routing}), bounded below
by the circuit depth $\Rmin$, and the compilation time; our implementation
additionally reports the average gate parallelism and the average non-routed-gate
percentage per layer (Definition~\ref{def:layering}).

\subsection{Results}

Table~\ref{tab:results} reports the routing-step count $|\Rout|$ and compilation
time across the suite for the Gaussian method and the two baselines, and
Figure~\ref{fig:scatter} plots per-circuit $|\Rout|$ of the Gaussian method against
WISQ: points below
the diagonal favour our method, while the WISQ column of the table flags circuits
where WISQ fails or times out though the Gaussian method completes. The field
weights of Section~\ref{sec:fields} were tuned per gaussian variant (coarse or fine) 
and safe-passage strategy (cube or connectivity) 
by benchmark sweeps, with the full values provided in the project repository~\cite{ftqc}.

% ---- TABLE: main comparison (single column) ----
\begin{table}[t]
    \centering
    \caption{Routing steps and compilation time per benchmark.\\[2pt]
    ``Gauss'' is the (tuned) Gaussian ``fine'' placement,\\
    ``Rand'' is the best of \todo{N} randomized baselines. \\
    $L_{\min}$ is the circuit layer count (minimum routing steps).\\
    $|\Q_L|$ is the circuit qubit count.\\
    \todo{Fill table from CSV runs + WISQ.}
    }
    \label{tab:results}
    \setlength{\tabcolsep}{3pt}
    \scriptsize
    \begin{tabular}{l r r rrr rrr}
        \toprule
        \textbf{Circuit} & \textbf{$|\Q_L|$} & \textbf{$L_{\min}$}
            & \multicolumn{3}{c}{\textbf{Routing steps}}
            & \multicolumn{3}{c}{\textbf{Compilation time (s)}} \\
        \cmidrule(lr){4-6}\cmidrule(lr){7-9}
            & & & \textbf{Gauss} & \textbf{Rand} & \textbf{WISQ}
              & \textbf{Gauss} & \textbf{Rand} & \textbf{WISQ} \\
        \midrule
        \texttt{adder}        & 28 & -- & -- & -- & -- & -- & -- & -- \\
        \texttt{qft}          & 20 & -- & -- & -- & -- & -- & -- & -- \\
        \texttt{multiplier}   & 45 & -- & -- & -- & -- & -- & -- & -- \\
        \texttt{square\_root} & 45 & -- & -- & -- & -- & -- & -- & -- \\
        \texttt{bwt}          & 21 & -- & -- & -- & -- & -- & -- & -- \\
        \texttt{hhl}          & 10 & -- & -- & -- & -- & -- & -- & -- \\
        \texttt{factor247}    & 15 & -- & -- & -- & -- & -- & -- & -- \\
        ~~$\vdots$ & & & & & & & & \\
        \midrule
        \textbf{Mean}    & -- & -- & -- & -- & -- & -- & -- & -- \\
        \bottomrule
    \end{tabular}
    \vspace{-0.2cm}
\end{table}

% ---- FIGURE: per-circuit scatter vs WISQ (single column) ----
\begin{figure}[t]
    \centering
    \figslot{1.6in}
    \caption{Per-circuit routing step ratios: Gaussian / WISQ~\cite{wisq}.
    Points below the diagonal indicate fewer routing steps for our method.
    \todo{Generate from CSVs.}}
    \label{fig:scatter}
    \vspace{-0.2cm}
\end{figure}

% ---- TABLE: tuned Gaussian parameters per regime (full width) ----
% Currently kept in the project README ("Tuned Gaussian mapping parameters").
% To show it in the paper, remove the \iffalse ... \fi wrapper and re-point the
% "tuned per regime" sentence in sections/evaluation.tex to Table~\ref{tab:params}.
\iffalse
\begin{table*}[t]
    \centering
    \caption{Tuned Gaussian parameters by regime (weighting variant $\times$
    lattice geometry). Common to all configurations: \texttt{external\_weight}
    $=0$, \texttt{base\_gaussian\_weight}$=1$, \texttt{cnot\_low}$=0$,
    \texttt{number\_of\_magic\_states}$=-1$ (auto), \texttt{center\_circle}
    magic placement with a border margin, and naive routing with smart
    $T$-routing. ``mapped'' is the repulsion-field weight and $\kappa$ the
    confidence of \eqref{eq:sigma}; \textsc{magic-high} depends on the border
    margin $b$ (percent), while all other weights do not.}
    \label{tab:params}
    \setlength{\tabcolsep}{5pt}
    \footnotesize
    \begin{tabular}{l cccc ccccc}
        \toprule
        & \multicolumn{4}{c}{\textbf{Border-independent weights}}
        & \multicolumn{5}{c}{\textbf{\textsc{magic-high} by border $b$ (\%)}} \\
        \cmidrule(lr){2-5}\cmidrule(lr){6-10}
        \textbf{Regime} & \textbf{$\kappa$} & \textbf{\textsc{cnot-high}} & \textbf{mapped} & \textbf{\textsc{magic-low}}
            & $b{=}0$ & $b{=}5$ & $b{=}10$ & $b{=}20$ & $b{=}30$ \\
        \midrule
        \regime{coarse}{cube} & 0.999999 & 1.5 & 0   & 0              & 0               & 0.2 & 0.2 & 0.2 & 1.6 \\
        \regime{coarse}{conn} & 0.99999  & 1.5 & 2.0 & 0              & 0               & 0.2 & 0.8 & 0.8 & 0.2 \\
        \regime{fine}{cube}   & 0.999999 & 0.5 & 0   & $1^{\dagger}$  & $20^{\ddagger}$ & 3   & 1.6 & 0   & 0   \\
        \regime{fine}{conn}   & 0.99999  & 0.5 & 1.5 & 0              & 3               & 1.6 & 0.4 & 0.4 & 0.2 \\
        \bottomrule
    \end{tabular}

    \vspace{2pt}
    {\footnotesize $^{\dagger}$\,\textsc{magic-low}$=1$ for border $b\le10\%$, else $0$ (fine\,/\,cube only).\quad
    $^{\ddagger}$\,Robust argmin is $20$, but the curve is flat above ${\sim}1.6$.}
\end{table*}
\fi


\todo{Discussion, 2--4 sentences once numbers are in: (i) where the field method
beats or loses to WISQ and why; (ii) how many randomized restarts the random
baseline needs to match field placement, if ever, and at what grid-area cost;
(iii) scalability.}
